% Generated draft: requires mathematical and source review.
\documentclass[11pt]{article}
\usepackage[margin=25mm]{geometry}
\usepackage{fontspec}
\setmainfont{Latin Modern Roman}
\usepackage{amsmath,amssymb,mathtools,unicode-math}
\setmathfont{Latin Modern Math}
\usepackage{enumitem,needspace,xurl,hyperref}
\setlength{\parindent}{0pt}
\setlength{\parskip}{5pt}
\setlength{\emergencystretch}{3em}
\newcommand{\N}{\mathbb{N}}
\newcommand{\Z}{\mathbb{Z}}
\newcommand{\Q}{\mathbb{Q}}
\newcommand{\R}{\mathbb{R}}
\newcommand{\C}{\mathbb{C}}
\newcommand{\F}{\mathbb{F}}
\newcommand{\A}{\mathbb{A}}
\newcommand{\PP}{\mathbb{P}}
\newcommand{\E}{\mathbb{E}}
\newcommand{\eps}{\varepsilon}
\newcommand{\dd}{\,\mathrm{d}}
\begin{document}
\section*{UnsolvedMath ID: 5300071}
\textbf{Source problem code:} catalogue entry
\subsection{Definitions and mathematical statement}
Let $f\in\mathbb C[z]$ have degree $d\geq1$, with all roots in the unit disk, and let $\alpha$ be a root of multiplicity $m\in\{1,\ldots,d\}$. Use the usual relaxed Newton map
\[
N_{h,f}(z)=z-h\frac{f(z)}{f'(z)},\qquad 0<h\leq m,
\]
viewed as a rational map after cancellation of removable factors; in particular $N_{h,f}(\alpha)=\alpha$. Let $A_h^*(\alpha)$ be the connected component containing $\alpha$ of the basin $\{z:N_{h,f}^{\circ n}(z)\to\alpha\}$. Define the common immediate-basin set
\[
B(f,\alpha)=\bigcap_{0<h\leq m}A_h^*(\alpha),\qquad C_R=\{z\in\mathbb C:|z|=R\}.
\]
Prove that there is an absolute constant $c>0$, independent of $f,d,\alpha,m,R$, such that for every such polynomial and root and every real $R\geq3$, the set $C_R\cap B(f,\alpha)$ contains arcs of total Euclidean arclength at least
\[
\frac{2\pi R}{cd}.
\]
A literal arc formulation is the existence of pairwise disjoint arcs $J_j\subseteq C_R\cap B(f,\alpha)$, finite or countably many, with $\sum_j\operatorname{length}(J_j)\geq2\pi R/(cd)$. The same points on these arcs must lie in the immediate basin for every real relaxation parameter in the whole interval $(0,m]$; separate large arcs depending on $h$ would not suffice. The supplied item does not define $N_{h,f}$ explicitly, so its normalization should be checked against the usual formula above.
\subsection{Short English statement}
For every root of a polynomial with roots in the unit disk, find uniformly large arcs on every circle of radius at least three that lie in its immediate basin for all relaxation parameters $0<h\leq m$ simultaneously. The arc fraction should be bounded below by an absolute constant times $1/d$.
\subsection{Sources}
\begin{itemize}
\item Scott Sutherland, 1992 problem list, item C2, as collected in Stony Brook, \emph{Open Problems in Dynamical Systems}, \url{https://www.math.stonybrook.edu/open-problems-dynamical-systems}; accessed 2026-07-29, with source-TeX and PDF-text extraction recorded in the provenance.
\item The supplied research summary mentions Smale's alpha theory, hybrid Newton methods, and large seeds as related context, but supplies no theorem citation establishing this bound.
\item Local registry \texttt{data/unsolvedmath/problems.json}, permanent UnsolvedMath ID $5300071$, source code catalogue entry, \emph{Uniform access to roots for relaxed Newton maps}. Dataset attribution: Ulam AI and the UnsolvedMath Contributors.
\end{itemize}
\end{document}
